% part: intuitionistic-logic
% chapter: propositions-as-types
% section: reduction

\documentclass[../../../include/open-logic-section]{subfiles}

\begin{document}

\olfileid{int}{pty}{red}

\olsection{Reduksi}

Ing !!{derivation}s dedhuksi alami, aturan !!a{introduction} sing
ditututi aturan !!a{elimination} iku dalan mubeng sing bisa ``direduksi''.
Contone, $\Intro{\land}$ sing ditututi $\Elim{\land}$ ``padha mbatalake'',
amarga konklusi $\Elim{\land}$ tansah salah siji konjung saka konjungsi
sing diasilake \Intro{\land}, lan premis $\Intro{\land}$ iku loro konjung
kasebut. Mula yen mbutuhake !!a{derivation} saka salah siji konjung,
kita bisa langsung nggunakake !!{derivation} kasebut tanpa ngenalake
konjungsi banjur ngeliminasi. Kita nduweni
\begin{gather*}
  \bottomAlignProof
  \AxiomC{}
  \RightLabel{$\delta_1$}
  \DeduceC{$!A_1$}
  \AxiomC{}
  \RightLabel{$\delta_2$}
  \DeduceC{$!A_2$}
  \RightLabel{$\Intro{\land}$}
  \BinaryInfC{$!A_1 \land !A_2$}
  \RightLabel{$\Elim{\land}_i$}
  \UnaryInfC{$!A_i$}
  \DisplayProof\bottomAlignProof
  \quad
  \redone
  \quad
  \AxiomC{}
  \RightLabel{$\delta_i$}
  \DeduceC{$!A_i$}
  \DisplayProof
\end{gather*}

Panyederhanaan iki uga menehi gagasan relasi reduksi sing padha tumrap
terma bukti sing cocog. Yen $N_1$ iku terma bukti kanggo~$\delta_1$ lan
$N_2$ kanggo~$\delta_2$, kita bisa ngarani yen terma bukti saka bukti
ing kiwa direduksi (kanthi siji langkah) dadi terma bukti saka bukti
ing tengen:
\[
  \proj{i}{\pair{N_1}{N_2}} \redone N_i
\]
Ing kalkulus lambda mawa tipe, iki aturan \emph{reduksi beta} kanggo
tipe prodhuk.

Ing dedhuksi alami, pasangan inferensi kaya \Intro{\land} sing
ditututi \Elim{\land}, yaiku pasangan sing bisa direduksi, diarani
\emph{cut}. Ing kalkulus lambda mawa tipe, terma ing kiwa $\redone$
diarani \emph{redex}, lan terma ing tengen diarani \emph{asil kontraksi}.
Ing kalkulus lambda tanpa tipe, mung $(\lambd[x][N])Q$ sing dianggep
redex. Ing kalkulus lambda mawa tipe, sintakse ditambahi terma sing
nggunakake $\pair{N}{M}$, $\proj{i}{N}$, $\inj[!A]{i}{N}$,
$\dcase{N}{x_1}{M_1}{x_2}{M_2}$, lan $\abort{!A}{N}$, kanthi
redex sing cocog.

Kanthi cara sing padha, kita nduweni reduksi kanggo disjungsi:
\begin{gather*}
  \bottomAlignProof
  \AxiomC{}
  \RightLabel{$\delta$}
  \DeduceC{$!A_i$}
  \RightLabel{$\Intro{\lor}$}
  \UnaryInfC{$!A_1 \lor !A_2$}
  \AxiomC{$\Discharge{!A_1}{x_1}$}
  \RightLabel{$\delta_1$}
  \DeduceC{$!C$}
  \AxiomC{$\Discharge{!A_2}{x_2}$}
  \RightLabel{$\delta_2$}
  \DeduceC{$!C$}
  \DischargeRule{$\Elim{\lor}$}{x_1,x_2}
  \TrinaryInfC{$!C$}
  \DisplayProof\bottomAlignProof
  \quad
  \redone
  \quad
  \AxiomC{}
  \RightLabel{$\delta$}
  \DeduceC{$!A_i$}
  \RightLabel{$\delta_i$}
  \DeduceC{$!C$}
  \DisplayProof
\end{gather*}
Iki cocog karo reduksi tumrap terma bukti:
\begin{gather*}
  \dcase{\inj[!A_{3-i}]{i}{M}}{x_1}{N_1}{x_2}{N_2} \redone \Subst{N_i}{M}{x_i}
\end{gather*}
kanthi $M$ minangka terma bukti saka~$\delta$, lan $N_i$ minangka
terma bukti saka~$\delta_i$. Iki aturan reduksi beta kanggo tipe gunggung.
Ing kene, $\Subst{N_i}{M}{x_i}$ tegese ngganti kabeh kedadeyan bebas
variabel~$x_i$ ing $N_i$ nganggo~$M$, kanthi ngganti jeneng variabel
kaiket dhisik yen perlu supaya variabel bebas ora kacekel.

Reduksi tipe fungsi cocog karo ngilangi dalan mubeng saka
$\Intro\lif$ sing ditututi $\Elim\lif$.
\begin{gather*}
  \bottomAlignProof
  \AxiomC{$\Discharge{!A}{x}$}
  \RightLabel{$\delta$}
  \DeduceC{$!B$}
  \DischargeRule{$\Intro{\lif}$}{x}
  \UnaryInfC{$!A \lif !B$}
  \AxiomC{}
  \RightLabel{$\delta'$}
  \DeduceC{$!A$}
  \RightLabel{$\Elim{\lif}$}
  \BinaryInfC{$!B$}
  \DisplayProof\bottomAlignProof
  \quad
  \redone
  \quad
  \AxiomC{}
  \RightLabel{$\delta'$}
  \DeduceC{$!A$}
  \RightLabel{$\delta$}
  \DeduceC{$!B$}
  \DisplayProof
\end{gather*}
Kanggo terma bukti, iki padha karo reduksi beta lumrah:
\begin{gather*}
  (\lambd[\typeof{x}{!A}][N]M)
  \redone \Subst{N}{M}{x}
\end{gather*}

Kajaba konversi reduksi tumrap !!{proof}s, kita bisa nimbang konversi
permutasi. Konversi iki ditrapake marang (sub)!!{proof}s sing dipungkasi
kanthi \Elim{\lor} banjur ditututi aturan !!{elimination} liyane,
contone:
\begin{multline*}
  \bottomAlignProof
  \AxiomC{}
  \RightLabel{$\delta$}
  \DeduceC{$!A_i$}
  \RightLabel{$\Intro{\lor}$}
  \UnaryInfC{$!A_1 \lor !A_2$}
  \AxiomC{$\Discharge{!A_1}{x_1}$}
  \RightLabel{$\delta_1$}
  \DeduceC{$!C_1 \land !C_2$}
  \AxiomC{$\Discharge{!A_2}{x_2}$}
  \RightLabel{$\delta_2$}
  \DeduceC{$!C_1 \land !C_2$}
  \DischargeRule{$\Elim{\lor}$}{x_1,x_2}
  \TrinaryInfC{$!C_1 \land !C_2$}
  \RightLabel{\Elim{\land}[i]}
  \UnaryInfC{$!C_i$}
  \DisplayProof\bottomAlignProof
  \quad
  \redone
  \\
  \AxiomC{}
  \RightLabel{$\delta$}
  \DeduceC{$!A_i$}
  \RightLabel{$\Intro{\lor}$}
  \UnaryInfC{$!A_1 \lor !A_2$}
  \AxiomC{$\Discharge{!A_1}{x_1}$}
  \RightLabel{$\delta_1$}
  \DeduceC{$!C_1 \land !C_2$}
  \RightLabel{\Elim{\land}[i]}
  \UnaryInfC{$!C_i$}
  \AxiomC{$\Discharge{!A_2}{x_2}$}
  \RightLabel{$\delta_2$}
  \DeduceC{$!C_1 \land !C_2$}
  \RightLabel{\Elim{\land}[i]}
  \UnaryInfC{$!C_i$}
  \DischargeRule{$\Elim{\lor}$}{x_1,x_2}
  \TrinaryInfC{$!C_i$}
  \DisplayProof
\end{multline*}
Yen terma bukti saka $\delta$, $\delta_1$, lan $\delta_2$ yaiku $M$,
$N_1$, lan $N_2$, aturan reduksi sing cocog kanggo terma bukti, yaiku
terma $\lambd$, awujud:
\[
  \proj{i}{\dcase{M}{x_1}{N_1}{x_2}{N_2}} \redone 
\dcase{M}{x_1}{\proj{i}{N_1}}{x_2}{\proj{i}{N_2}}
\]

Kita bisa ngilangi dalan mubeng ora mung ing pungkasan !!{proof}s,
nanging uga ing sajrone !!{proof}s, kanthi ngetrapake konversi reduksi
marang sub-!!{proof} sing dipungkasi dalan mubeng. Semono uga, reduksi
$\beta$ ditrapake marang subterma saka terma $\lambd$. Yen $N$ iku
subterma saka $M$ lan $N \redone N'$, kita bisa ngganti subterma $N$
ing~$M$ nganggo $N'$ kanggo ngasilake $M'$. Ing kasus iki kita uga
nulis $M \redone M'$. Yen ana urutan $M = M_1 \redone M_2 \redone
\dots \redone M_k = M'$ (kalebu $k=1$, yaiku $M = M'$), kita nulis
$M \red M'$.

!!{proof} sing ora ngidini konversi ditrapake marang sub-!!{proof}
endi wae diarani \emph{normal}. Semono uga, terma $\lambd$ sing ora
ngidini konversi ditrapake (marang subterma endi wae) diarani
\emph{normal}. Terma $\lambd$ normal uga diarani \emph{wujud normal}.

\end{document}
